\documentclass[10pt]{article}

\usepackage[utf8]{inputenc}

\usepackage[T1]{fontenc}

\usepackage{amsmath}

\usepackage{amsfonts}

\usepackage{amssymb}

\usepackage[version=4]{mhchem}

\usepackage{stmaryrd}

\usepackage{mathrsfs}

\usepackage{bbold}

\usepackage{graphicx}

\usepackage[export]{adjustbox}

\graphicspath{ {./images/} }



\begin{document}

СПЕКТРАЛЬНЫЙ ОПЕРАТОР - ограпиченный линейный ошератор $A$, отображающий банахово пространство $X$ в себя и такой, что для $\sigma$-алгебры $\mathscr{B}$ борелевских множеств $\delta$ на плоскости существует разложение единицы $E(\delta)$ со свойствами: 1) для любого $\delta \in \mathscr{B}$ проектор $E(\delta)$ приводит $A$, т. е. $E(\delta) A=A E(\delta)$, и спектр $\sigma\left(A_{\delta}\right)$ лежит в $\bar{\delta}$, где $A_{\delta}$ - сужение оператора $A$ на инвариантное подпространство $E(\delta) X ; 2)$ отображение $\delta \rightarrow E(\delta)$ есть гомоморфизм $\mathscr{B}=\{\delta\}$ в булеву алгебру $\{E(\delta)\} ; 3)$ все проекторы $E(\delta)$ ограничены, т. е. $\|E(\delta)\| \leqslant$ $\leqslant M, \delta \in \mathscr{B} ; 4)$ разложение единицы $E(\delta)$ счетно аддитивно в сильной топологии пространства $X$, т. е. для любого $x \in X$ и любой последовательности $\left\{\delta_{n}\right\} \subset \mathscr{B}$, состоящей из попарно непересекающихся множеств,



\[

E\left(\cup_{n=1}^{\infty} \delta_{n}\right) x=\sum_{n=1}^{\infty} E\left(\delta_{n}\right) x .

\]



Понятие С. о. можно распространить на неограниченные замкнутые операторы. При этом в 1) надо дополиительно потребовать, чтобы выполнялось включение $E(\delta) D(A) \subset D(A)$, где $D(A)-$ область определения оператора $A$, и $E(\delta) X \subset D(A)$ для ограниченных $\delta$.



С. о. являются все линейные операторы в конечномерном пространстве, самосопряженные и нормальные операторы в гильбертовом простравстве, напр. оператор



\[

A x(t)=t x+\int_{-\infty}^{\infty} K(t, s) x(s) d s

\]



в $L_{p}(-\infty, \infty), 1<p<\infty$, на



\[

D(A)=\left\{\left.x(t)\left|\int_{-\infty}^{\infty}\right| t x(t)\right|^{2} d t<\infty,\right.

\]



если ядро $K(t, s)$ есть преобразование Фурье борелев. ской меры $\mu$ на плоскости с полной вариацией var $\mu<1 / 2 \pi$ и такое, что



\[

\int_{-\infty}^{\infty} K(t, s) x(s) d s, \int_{-\infty}^{\infty} K(t, s) x(t) d t

\]



суть ограниченные линейные операторы в $L_{p}(-\infty, \infty)$.



С. о. обладают рядом важных свойств, напр.:



\[

\lambda \in \delta(A) \Leftrightarrow \exists\left\{x_{n}\right\} \subset X,\left\|x_{n}\right\|=1,(A-\lambda I) x_{n} \rightarrow 0

\]



в случае сегіарабельного $X$ тотечный и остаточный спектры $A$ не более чем счетны и др.



Лum.: [1] Д а н ф о р д Н., ІІІ в а р ц д ж., Линейныс опе[2] Д, а н ф о р д Н., «Математика», 1960, т. 4, в. 1, с. 53-100.



СПЕКТРАЛЬНЫЙ РАДИУС э л е м е т т а б в $а$ х о в ой а л г е б ры - раднус $\rho$ наименьшего круга на плоскости, содержащего спектр этого әлемснта. С. p. элемента a связан с нормами его с'епеней формулой



\[

\rho(a)=\lim _{n \rightarrow \infty}\left\|a^{n}\right\|^{1 / n}=\inf \left\|a^{n}\right\|^{1 / n}

\]



из к-рой следует, в частности, что $\rho(a) \leqslant\|a\|_{\text {. }}$ С. р. ограниченного оператора в банаховом пространстве это его C. р. как элемента банаховой алгебры всех операторов. В гильбертовом пространстве С. р. оператора равен точной нижней грани норм подобных ему операторов (см. [2]):



\[

\rho(A)=\inf _{X}\left\|X A X^{-1}\right\| .

\]



Если оператор нормален, то $\rho(A)=\|A\|$.



С. p.- полунепрерывная сверху (но, вообще говоря, не непрерывная) функция әлемента банаховой алгебры. Доказана [3] субгармоничность С. р. (это означает, что если $z \rightarrow h(z)$ - голоморфное отображение нек-рой облласти $D \subset \mathbb{C}$ в банахову алгебру $\mathfrak{A}$, то $z \rightarrow \rho(h(z))$ - субгармонич. функция). Лит.: [1] H а й м а р к М. А., Нормированные кольца,

изд., М., 1968; [2] Х а л м о ш П., Гильбертово пространство 2 изд., М., 192, с англ., M., 1970 ; [3] V e s e n t i n i E., "Boll. London Math. Soc.", 1970, v. 2, p. 327-34. B. С. IIульман.



СПЕКТРАЛЬНЫЙ СЕМИИНВАРИАНТ - одна из характеристик стационарного случайного процесса. Пусть $X(t),-\infty<t<\infty,-$ действительвый стационарный случайный процесс, для к-рого $E|X(t)|^{n} \leqslant$ $\leqslant C<\infty$. Семивнварианты әтого процесса



\[

\begin{gathered}

S^{(n)}\left(t_{1} \ldots t_{n}\right)= \\

=\left.\frac{i^{-n} \partial^{n}}{\partial u_{1} \ldots \partial u_{n}} \ln \mathrm{E} e^{i\left(u_{1} X\left(t_{1}\right)+\ldots+u_{n} X\left(t_{n}\right)\right)}\right|_{u_{1}=\ldots u_{n}=0}

\end{gathered}

\]



связаны с моментами



\[

M^{(n)}\left(t_{1} \ldots t_{n}\right)=\mathrm{E}\left\{X\left(t_{1}\right) \ldots X\left(t_{n}\right)\right\}

\]



соотнотениями



$S^{(n)}(I)=\sum_{\mathrm{U}_{p=1}^{q} I_{p}=I}(-1)^{q-1}(q-1) ! \prod_{p=1}^{q} M^{(p)}\left(I_{p}\right)$



\[

M^{(n)}(I)=\sum_{\mathrm{U}_{p=1}^{q} I_{p}=I} \prod_{p=1}^{q} S^{(p)}\left(I_{p}\right),

\]



где



\[

I=\left(t_{1} \ldots t_{n}\right), I_{p}=\left(t_{i_{1}} \ldots t_{i_{p}}\right) \subseteq I

\]



и суммирование ведется по всем разбиениям множества $I$ на непересекающиеся подмножества $I_{p}$. Говорят, что $X(t) \in \Phi^{(n)}$, если для всех $1 \leqslant k \leqslant n$ в пространстве $\mathbb{R}^{k}$ существует мера $M^{(k)}(\Delta)$ ограниченной вариации такая, что для всех $t_{1} \ldots t_{k}$



$M^{(k)}\left(t_{1} \ldots t_{k}\right)=\int_{\mathbb{R}^{k}} e^{i\left(t_{1} \lambda_{1}+\ldots+t_{k} \lambda_{k}\right)} M^{(k)}\left(d \lambda_{1} \ldots d \lambda_{k}\right)=$



\[

=\int_{\mathbb{R}^{k}} e^{i(t, \lambda)} M^{(k)}(d \lambda)

\]



Меру $F^{(n)}$, оиределенную на системе борелевских множеств, наз. с пі е к т р аль ным с е ми и н в а р иа н т о м, если для всех $t_{1} \ldots t_{n}$



\[

S^{(n)}\left(t_{1} \ldots t_{n}\right)=\int_{\mathbb{R}^{n}} e^{i(t, \lambda)} F^{(n)}(d \lambda) .

\]



Мера $F^{(n)}$ существует и имеет огрангченную вариацию, если $X(t) \in \Phi^{(n)}$. В случае стационарного процесса $X(t)$ семиинварианты $S^{(n)}\left(t_{1} \ldots t_{n}\right) \quad$ инвариантны относительно сдвигов:



\[

S^{(n)}\left(t_{1}+\tau, \ldots, t_{n}+\tau\right)=S^{(n)}\left(t_{1} \ldots t_{n}\right),

\]



a спектральные меры $F^{(n)}$ и $M^{(n)}$ сосредоточены на многообразии $\lambda_{1}+\ldots+\lambda_{n}=0$. Если мера $F^{(n)}$ абсолютно непрерывна относительно меры Лебега на этом многообразии, то существует спектральная плотность $n$-го порядка $f_{n}\left(\lambda_{1}, \ldots, \lambda_{n-1}\right)$, определяемая равенствами



\[

\begin{aligned}

S^{(n)}\left(t_{1} \ldots t_{n}\right)= & \int_{\mathbb{R}^{n-1}} e^{i\left(\lambda_{1}\left(t_{2}-t_{1}\right)+\ldots+\lambda_{n-1}\left(t_{n}-t_{1}\right)\right)} \times \\

& \times f_{n}\left(\lambda_{1}, \ldots, \lambda_{n-1}\right) d \lambda,

\end{aligned}

\]



верными при всех $t_{1} \ldots t_{n}$. В случае дискретного времени под $\mathbb{K}^{(k)}$ во всех приведенных выше формулах надо понимать $k$-мерный куб $-\pi \leqslant \lambda_{1} \leqslant \pi, 1 \leqslant i \leqslant k$.



\includegraphics[max width=\textwidth, center]{2023_07_09_1b60e44b844966e8feddg-1}

рия вероятностей, 2 изд., М., $1973 ;$ [2] J е о н о в В. П., Не-

которые применения старших семиинвариантов к теории стакоторые применения старших семиинвариантов к теории ста-

ционарных случайных процессов, М., 1964. И. Г. Журбенко.



СПЕКТРАЛЬНЫІИ СИНТЕЗ - восстановление инвариантных подпространств семейства линейных операторов по содержащимся в них собственным или корневым шодпространствам этого семейства. Точнее, пусть $\mathfrak{X}$ - коммутативное семейство операторов в то ІІологическом векторном пространстве $X, \quad v_{p}(\mathfrak{A}) \cdots$ его точеqный спектр, т. е. совокупность числовых фұфук-





\end{document}